\section{Introduction}\label{sec:introduction}

For a finite loopless undirected multigraph \(G=(V,E)\) and a nonempty proper
subset \(S\subset V\), write \(\delta_G(S)\) for the edges with exactly
one endpoint in \(S\). Parallel edges are counted separately.
A graph is \(k\)-edge-connected if \(|\delta_G(S)|\ge k\) for
every nonempty proper \(S\).
A spanning tree \(T\) is \(\alpha\)-thin if
\[
 |\delta_T(S)|\le\alpha|\delta_G(S)|
 \qquad(\varnothing\ne S\subsetneq V).
\]
A spanning tree retains connectivity using only \(|V|-1\) edges.
Thinness asks that these edges occupy a small fraction of every cut
simultaneously. The question is whether high edge connectivity alone
forces such a tree, with a bound independent of the number of vertices.

\begin{theorem}\label{thm:main}
There is a universal constant \(C>0\) such that, for every integer
\(k\ge1\), every finite \(k\)-edge-connected undirected multigraph
without loops and with at least two vertices has a spanning tree \(T\)
satisfying
\[
 |\delta_T(S)|\le\frac Ck\,|\delta_G(S)|
 \qquad\text{for every }\varnothing\ne S\subsetneq V.
\]
\end{theorem}

\subsection*{The conjecture and its history}

Goddyn's thin tree conjecture asks whether, for every
\(\varepsilon>0\), a sufficiently high edge connectivity guarantees
an \(\varepsilon\)-thin spanning tree~\cite{Goddyn}.
The strong thin tree conjecture asks for the quantitative rate \(C/k\).
Theorem~\ref{thm:main} resolves this stronger conjecture positively.
The order \(1/k\) is necessary: with two vertices and \(k\) parallel
edges, every spanning tree uses exactly one of the \(k\) edges of
the only cut.

Thin trees connect this graph-theoretic question with approximation
algorithms. In the asymmetric traveling salesman problem, one seeks a
minimum-cost tour through all vertices of a directed metric.
Asadpour, Goemans, M\k{a}dry, Oveis Gharan, and Saberi used
maximum-entropy distributions on spanning trees to round the standard
linear-programming relaxation, obtaining an
\(O(\log n/\log\log n)\)-approximation, where \(n\) is the number
of vertices~\cite{AsadpourSODA,Asadpour}.
Their method explains the usefulness of controlling every cut:
together with control of the tree's cost, cut bounds permit an inexpensive
Eulerian augmentation of a suitable oriented tree.
For the underlying thin-tree problem, Oveis Gharan and
Saberi established the inverse-connectivity bound for planar graphs
and, with constants depending on the surface, for graphs of bounded
orientable genus~\cite[Theorem~5.1]{OveisGharanSaberi2011}.

For general multigraphs, Anari and Oveis Gharan developed
effective-resistance-reducing flows and the geometry of locally
connected hierarchies. They proved the existence of a spanning tree
with thinness \(\operatorname{poly}(\log\log n)/k\)
\cite[Corollary~1.8]{AOG}.
Their work provides the principal geometric foundation of the present
proof. Progress for specified families of cuts gives a complementary
view of the problem. Klein and Olver obtained inverse-connectivity
thinness for every prescribed laminar family~\cite{KleinOlver2023}.
Klein, Olver, and Yeoh constructed a spanning tree crossing every cut
of size less than \((1+1/40)k\) at most \(88\) times
\cite[Theorem~1]{KleinOlverYeoh}.
Theorem~\ref{thm:main} controls all cuts simultaneously and removes
the dependence on \(n\).

The thin-tree approach to asymmetric traveling salesman remains
structurally useful, although constant-factor approximation for that
problem has already been obtained by other methods. Svensson,
Tarnawski, and V\'egh proved the first such guarantee
\cite{SvenssonTarnawskiVegh2020}; Traub and Vygen subsequently
improved it~\cite{TraubVygen2022}. Our objective here is the
underlying all-cut spanning-tree theorem.

\subsection*{The two main deductions}

High edge connectivity supplies many edge-disjoint spanning trees
through the Nash--Williams--Tutte theorem~\cite{NashWilliams,Tutte}.
For any fixed cut, the average number of edges used by the trees in
such a packing is small. This averaging does not select one tree
that is small on every cut. We maintain a large packing throughout
a sequence of sparsifications: each step reduces cut sizes and the
packing size by almost the same factor, and the accumulated relative
loss remains bounded.

The first deduction extracts a large tree packing on edges of small
resistance. For a graph \(H\), let \(b_e=\mathbf1_u-\mathbf1_v\)
be an arbitrarily oriented incidence column of an edge \(e=uv\),
and let \(L_H=\sum_e b_eb_e^{\mathsf T}\) be its Laplacian.
The resistance of \(e\) with respect to a positive definite matrix
\(X\) is \(b_e^{\mathsf T}X^{-1}b_e\).
Write \(\mathbf1_S\) for the indicator column of \(S\), and
\(A\preceq B\) when \(B-A\) is positive semidefinite.
For all sufficiently large integers \(r\), if \(H\) on at least
two vertices contains \(r\) edge-disjoint spanning trees, we
construct such a matrix \(X\) with
\[
 L_H\preceq X,\qquad
 \mathbf1_S^{\mathsf T}X\mathbf1_S
 \le3|\delta_H(S)|
 \quad(\varnothing\ne S\subsetneq V),
\]
such that the edges of \(H\) with
\(b_e^{\mathsf T}X^{-1}b_e\le4r^{-1/2}\) still contain
\(\lfloor(1-r^{-1/16})r\rfloor\) edge-disjoint spanning trees.
The starting point is a polynomial consequence of the locally connected
hierarchy theorem of Anari and Oveis Gharan~\cite[Theorem~3.5]{AOG};
we prove the required consequence in full.
For a given matrix, retain the edges below each resistance threshold
and partition the vertices into maximal sets supporting the prescribed
number of disjoint spanning trees. These partitions are nested.
Selecting them by dyadic crossing-edge counts gives a rooted hierarchy
of nested vertex sets and avoids a loss depending on the number of
levels. The hierarchy estimate supplies a response matrix \(Y\) with
small resistance on almost all crossing edges of every selected
partition at once (Lemma~\ref{lem:response}). We interpolate finitely
many such responses and take a fixed-point limit. At the resulting
matrix \(X\), a nontrivial maximal packing partition of its own
low-resistance graph would have too many crossing edges. This
contradiction yields the extracted packing
(Theorem~\ref{thm:extraction}).

The second deduction is a sparsification step that retains about half
the tree packing while bounding each new cut by approximately half its
previous size. We use two vector representations of the incidence
columns. The matrix \(X\) controls cut size. A separate isotropic
representation---whose vector outer products sum to the identity---
controls the partition inequalities for tree packing.
Applying the Marcus--Spielman--Srivastava theorem~\cite{MSS} to the stacked
columns controls both systems on the same selected edge set.
The relative losses are \(1+O(r^{-1/16})\). They have a bounded product
as the packing parameter decreases geometrically.
Harvey and Olver used iterated matrix halving to obtain spectrally thin
trees under an effective-resistance hypothesis
\cite[Appendix~F]{HarveyOlver2014}. The additional packing
representation here permits iteration from edge connectivity.

Preserving a packing requires more than a weighted spectral
approximation. Sampling by effective resistance and spectral
sparsification provide powerful ways to reduce a graph while
approximately preserving its Laplacian
\cite{SpielmanSrivastava2011,BatsonSpielmanSrivastava2012}.
Here the edges must retain their original unit multiplicities.
For example, \(q\) parallel edges between two vertices have the same
Laplacian as one edge of weight \(q\), but the former graph packs
\(q\) edge-disjoint spanning trees and the latter has only one.
The second representation supplies the rank information needed to
certify the partition inequalities after selecting an unweighted
subset of edges.

\subsection*{Organization and inputs}

Section~\ref{sec:preliminaries} proves the tree-packing facts used below.
Section~\ref{sec:hierarchy} gives the precise hierarchy input, and
Section~\ref{sec:extraction} proves the extraction theorem.
Sections~\ref{sec:auxiliary} and~\ref{sec:sparsification} develop the
auxiliary representation and the sparsification step.
Section~\ref{sec:iteration} proves Theorem~\ref{thm:main}, including
the small-connectivity cases.
The independent-vector theorem used in the sparsification appears
as Theorem~\ref{thm:mss}. The required hierarchy estimate and the
finite-dimensional fixed-point theorem are proved in the appendices.
All sufficiently large thresholds in the argument are absolute.
Choices made after a particular finite graph has been fixed need
not have a uniform algorithmic running time.
The companion article \emph{A polynomial-time construction of strong
thin trees}~\cite{ConstructiveThinTreesCompanion} develops the
separate finite-bit construction, including graphs given by
binary-encoded parallel-edge multiplicities. Its geometric input
extends the hierarchy proof here from cut tests to a larger cone of
metric tests.
